While we could use the Sieve of Eratosthenes again, this problem is quite famous for Lucy_Hedgehogs' legendary forum post.
We analyse it here as this is a pretty important algorithm and worth a thorough analysis now. We quote their analysis here:

<blockquote>
Let $S(v, m)$ be the sum of integers in the range that remain after sieving with all primes smaller or equal than $m$.
That is, $S(v, m)$ is the sum of integers up to $v$ that are either prime or the product of primes larger than $m$.

$S(v, p)$ is equal to $S(v, p-1)$ if $p$ is not prime or $v$ is smaller than $p^2$. Otherwise $S(v, p)$ can be computed
from $S(v, p-1)$ by finding the sum of integers that are removed while sieving with $p$. An integer is removed in this
step if it is the product of $p$ with another integer that has no divisor smaller than $p$. This can be expressed as

$$S(v, p)=S(v, p-1)-p\left(S\left(\frac{v}{p}, p-1\right)-S(p-1, p-1)\right)$$

Dynamic programming can be used to implement this. It is sufficient to compute $S(v, p)$ for all positive integers $n$
that are representable as $\lfloor n/k\rfloor$ for some integer $k$ and all $p\le v$.
</blockquote>